% GENERATED FILE. Edit canonical lessons, then run npm run generate:books.
% canonical-source-hash: fnv1a64:c1dede0357388629
% canonical-lessons: SA-C337-vanavasah, SA-C337-pratijna, SA-C337-paduka, SA-C337-pratijanite, SA-C337-mrgaya

\chapter{Exile in the Forest}
\label{ch:sa-exile-in-the-forest}
\hlchaptermodality{337}

\begin{chapteropening}
\textbf{By the end of this chapter:} \emph{I can tell how a prince kept his word in exile.}

Complete the last lesson of chapter 337: I can tell how a prince kept his word in exile.
\end{chapteropening}

\section[\texorpdfstring{\sk{वनवासः} (vanavāsaḥ)}{वनवासः (vanavāsaḥ)}]{\sk{वनवासः} (vanavāsaḥ) — life in the forest; exile in the forest}

\label{lesson:SA-C337-vanavasah}



\begin{quote}
Before the new one: say the Sanskrit for bewildered, then the Sanskrit for tempts.
\end{quote}

\subsection*{You'll want to know: \sk{वनवासः}}
\textbf{\sk{वनवासः}} — \emph{vanavāsaḥ} — \textquotedblleft{}life in the forest; exile in the forest\textquotedblright{}.

\begin{quote}
\sk{वनवासे} \sk{सः} \sk{फलानि} \sk{खादति।}
\end{quote}

\begin{quote}
\emph{vanavāse saḥ phalāni khādati.}
\end{quote}

\begin{quote}
In his forest exile he lives on fruit.
\end{quote}

A masculine noun.

\begin{cousinweb}
\textbf{\sk{वनवासः}} (\emph{vanavāsaḥ}) is built from \emph{vana}, \textquotedblleft{}a forest\textquotedblright{}, and \emph{vāsa}, \textquotedblleft{}dwelling\textquotedblright{}, from the root √vas, as in \emph{vasati}; MW gives \emph{vanavāsa} as dwelling in a forest (MW).
\end{cousinweb}

\begin{grammarlens}[title={What you've built}]
One more word for telling what happened.
\end{grammarlens}

\subsection*{Your turn}
\begin{itemize}
\raggedright
  \item \emph{Say it:} \emph{vanavāsaḥ}
  \item \emph{Say it:} \emph{vanavāsaḥ}, once more
  \item \emph{Say it:} \emph{vanavāse saḥ phalāni khādati.} — In his forest exile he lives on fruit.
  \item \emph{Recall:} say the Sanskrit for bewildered, then the Sanskrit for tempts, then say \emph{vanavāsaḥ} again
\end{itemize}

\begin{tcolorbox}[breakable,colback=teal!4,colframe=teal!35!black,title={Before you move on}]
What does \sk{वनवासः} mean? (\textquotedblleft{}Life in the forest; exile in the forest\textquotedblright{}.) Say it once more.
\end{tcolorbox}
\section[\texorpdfstring{\sk{प्रतिज्ञा} (pratijñā)}{प्रतिज्ञा (pratijñā)}]{\sk{प्रतिज्ञा} (pratijñā) — a promise; a vow}

\label{lesson:SA-C337-pratijna}



\begin{quote}
Before the new one: say the Sanskrit for tempts, then the Sanskrit for life in the forest.
\end{quote}

\subsection*{You'll want to know: \sk{प्रतिज्ञा}}
\textbf{\sk{प्रतिज्ञा}} — \emph{pratijñā} — \textquotedblleft{}a promise; a vow\textquotedblright{}.

\begin{quote}
\sk{सः} \sk{पितुः} \sk{प्रतिज्ञां} \sk{पालयति।}
\end{quote}

\begin{quote}
\emph{saḥ pituḥ pratijñāṃ pālayati.}
\end{quote}

\begin{quote}
He keeps his father's promise.
\end{quote}

A feminine noun.

\begin{cousinweb}
\textbf{\sk{प्रतिज्ञा}} (\emph{pratijñā}) is built from \emph{prati} and the root √jñā, \textquotedblleft{}to know\textquotedblright{}; MW gives \emph{pratijñā} as admission, agreement, promise, vow (MW).
\end{cousinweb}

\begin{grammarlens}[title={What you've built}]
One more word for telling what happened.
\end{grammarlens}

\subsection*{Your turn}
\begin{itemize}
\raggedright
  \item \emph{Say it:} \emph{pratijñā}
  \item \emph{Say it:} \emph{pratijñā}, once more
  \item \emph{Say it:} \emph{saḥ pituḥ pratijñāṃ pālayati.} — He keeps his father's promise.
  \item \emph{Recall:} say the Sanskrit for tempts, then the Sanskrit for life in the forest, then say \emph{pratijñā} again
\end{itemize}

\begin{tcolorbox}[breakable,colback=teal!4,colframe=teal!35!black,title={Before you move on}]
What does \sk{प्रतिज्ञा} mean? (\textquotedblleft{}A promise; a vow\textquotedblright{}.) Say it once more.
\end{tcolorbox}
\section[\texorpdfstring{\sk{पादुका} (pādukā)}{पादुका (pādukā)}]{\sk{पादुका} (pādukā) — a sandal}

\label{lesson:SA-C337-paduka}



\begin{quote}
Before the new one: say the Sanskrit for life in the forest, then the Sanskrit for a promise.
\end{quote}

\subsection*{You'll want to know: \sk{पादुका}}
\textbf{\sk{पादुका}} — \emph{pādukā} — \textquotedblleft{}a sandal\textquotedblright{}.

\begin{quote}
\sk{भ्राता} \sk{तस्य} \sk{पादुके} \sk{सिंहासने} \sk{स्थापयति।}
\end{quote}

\begin{quote}
\emph{bhrātā tasya pāduke siṃhāsane sthāpayati.}
\end{quote}

\begin{quote}
His brother places his sandals on the throne.
\end{quote}

A feminine noun.

\begin{cousinweb}
\textbf{\sk{पादुका}} (\emph{pādukā}) is usually traced to \emph{pāda}, \textquotedblleft{}a foot\textquotedblright{}; MW gives \emph{pādukā} as a shoe or slipper, and the footprint of a holy person (MW).
\end{cousinweb}

\begin{grammarlens}[title={What you've built}]
One more word for telling what happened.
\end{grammarlens}

\subsection*{Your turn}
\begin{itemize}
\raggedright
  \item \emph{Say it:} \emph{pādukā}
  \item \emph{Say it:} \emph{pādukā}, once more
  \item \emph{Say it:} \emph{bhrātā tasya pāduke siṃhāsane sthāpayati.} — His brother places his sandals on the throne.
  \item \emph{Recall:} say the Sanskrit for life in the forest, then the Sanskrit for a promise, then say \emph{pādukā} again
\end{itemize}

\begin{tcolorbox}[breakable,colback=teal!4,colframe=teal!35!black,title={Before you move on}]
What does \sk{पादुका} mean? (\textquotedblleft{}A sandal\textquotedblright{}.) Say it once more.
\end{tcolorbox}
\section[\texorpdfstring{\sk{प्रतिजानीते} (pratijānīte)}{प्रतिजानीते (pratijānīte)}]{\sk{प्रतिजानीते} (pratijānīte) — promises}

\label{lesson:SA-C337-pratijanite}



\begin{quote}
Before the new one: say the Sanskrit for a promise, then the Sanskrit for a sandal.
\end{quote}

\subsection*{You'll want to know: \sk{प्रतिजानीते}}
\textbf{\sk{प्रतिजानीते}} — \emph{pratijānīte} — \textquotedblleft{}promises\textquotedblright{}.

\begin{quote}
\sk{सः} \sk{पित्रे} \sk{प्रतिजानीते}, \sk{अहं} \sk{वनं} \sk{गमिष्यामि} \sk{इति।}
\end{quote}

\begin{quote}
\emph{saḥ pitre pratijānīte, ahaṃ vanaṃ gamiṣyāmi iti.}
\end{quote}

\begin{quote}
He promises his father that he will go to the forest.
\end{quote}

Past (he …): \sk{सः} \sk{पित्रे} \sk{प्रतिज्ञातवान्}, \emph{saḥ pitre pratijñātavān}.

\begin{cousinweb}
\textbf{\sk{प्रतिजानीते}} (\emph{pratijānīte}) is built from \emph{prati} and the root √jñā, \textquotedblleft{}to know\textquotedblright{}; MW gives \emph{prati-jñā} as to admit, acknowledge, and to promise (MW).
\end{cousinweb}

\begin{grammarlens}[title={What you've built}]
One more word for telling what happened.
\end{grammarlens}

\subsection*{Your turn}
\begin{itemize}
\raggedright
  \item \emph{Say it:} \emph{pratijānīte}
  \item \emph{Say it:} \emph{pratijānīte}, once more
  \item \emph{Say it:} \emph{saḥ pitre pratijānīte, ahaṃ vanaṃ gamiṣyāmi iti.} — He promises his father that he will go to the forest.
  \item \emph{Recall:} say the Sanskrit for a promise, then the Sanskrit for a sandal, then say \emph{pratijānīte} again
\end{itemize}

\begin{tcolorbox}[breakable,colback=teal!4,colframe=teal!35!black,title={Before you move on}]
What does \sk{प्रतिजानीते} mean? (\textquotedblleft{}Promises\textquotedblright{}.) Say it once more.
\end{tcolorbox}
\section[\texorpdfstring{\sk{मृगया} (mṛgayā)}{मृगया (mṛgayā)}]{\sk{मृगया} (mṛgayā) — hunting; the chase}

\label{lesson:SA-C337-mrgaya}



\begin{quote}
Before the new one: say the Sanskrit for a sandal, then the Sanskrit for promises.
\end{quote}

\subsection*{You'll want to know: \sk{मृगया}}
\textbf{\sk{मृगया}} — \emph{mṛgayā} — \textquotedblleft{}hunting; the chase\textquotedblright{}.

\begin{quote}
\sk{राजा} \sk{मृगयायै} \sk{वनं} \sk{गच्छति।}
\end{quote}

\begin{quote}
\emph{rājā mṛgayāyai vanaṃ gacchati.}
\end{quote}

\begin{quote}
The king goes to the forest to hunt.
\end{quote}

A feminine noun.

\begin{cousinweb}
\textbf{\sk{मृगया}} (\emph{mṛgayā}) is usually traced to \emph{mṛga}, \textquotedblleft{}a wild animal, a deer\textquotedblright{}; MW gives \emph{mṛgayā} as hunting, the chase (MW).
\end{cousinweb}

\begin{grammarlens}[title={What you've built}]
One more word for telling what happened.
\end{grammarlens}

\subsection*{Your turn}
\begin{itemize}
\raggedright
  \item \emph{Say it:} \emph{mṛgayā}
  \item \emph{Say it:} \emph{mṛgayā}, once more
  \item \emph{Say it:} \emph{rājā mṛgayāyai vanaṃ gacchati.} — The king goes to the forest to hunt.
  \item \emph{Recall:} say the Sanskrit for a sandal, then the Sanskrit for promises, then say \emph{mṛgayā} again
\end{itemize}

\begin{tcolorbox}[breakable,colback=teal!4,colframe=teal!35!black,title={Before you move on}]
What does \sk{मृगया} mean? (\textquotedblleft{}Hunting; the chase\textquotedblright{}.) Say it once more.
\end{tcolorbox}
